\section{Symmetric counts and their domination}\label{sec:counts}
Remove empty requirements, merge $\rho$-orbits, and discard an orbit
contained in a different larger orbit. If $s$ is contained in $t$, then
$\rho(s)$ is contained in $\rho(t)$, so either orientation of the larger
word satisfies an orientation of the smaller. This preprocessing preserves
feasibility and the optimum. If nothing remains, return $\eps$.
Otherwise let $R$ include both orientations of the surviving requirements,
and let $V$ consist of all their substrings, including $\eps$.
It is substring-closed and $\rho$-closed; required words are maximal in $V$.
Let $L$ be the total symbol length of one representative per surviving
orbit, so $L\le N$ and $|V|=O(L^2)$.

\subsection{Blocking and the recursion}\label{sec:blocking}
A periodic text is a map $A:\ZZ\to\Sigma$ with $A(i+p)=A(i)$ for
some positive $p$, not necessarily least. For $s=A[0:h)$ and $r\in V$,
let $b_{A,s}(r)$ count the integers $j$ such that
\begin{gather}
1\le j,\quad j+h\le |r|-1,\quad r[j:j+h)=s,\notag\\
r(i)=A(i-j)\quad(1\le i\le|r|-2),\qquad
r(0)\ne A(-j),\quad r(|r|-1)\ne A(|r|-1-j).\label{eq:block}
\end{gather}
Put $B_A(s;m)=\sum_{r\in V\setminus\{\eps\}}b_{A,s}(r)m(r)$.
Only $|r|\ge |s|+2$ contributes. For an occurrence aligned with $A$
in a finite word, the nearest mismatches on both sides, if present,
determine a unique bracket. Call it blocked when that bracket belongs
to $V$; absence of either mismatch means unblocked. Consequently
$B_A(s;N_T)$ counts blocked occurrences exactly. This is the blocking
contract of \src{Section 2}; its coefficients count placements,
not just distinct bracket words.

Define $m(s)$ in decreasing length, simultaneously on each $\rho$-orbit.
Missing words have count zero. The lower bounds are
\begin{equation}\label{eq:extension}
m(s)\ge\sum_c m(cs),\qquad m(s)\ge\sum_c m(sc),
\end{equation}
with required lower bound $1$ for nonfixed words and $2$ for fixed words.
For $v=A[0:h)$, $w=A[0:h-kp)$ nonempty, $k\ge1$, impose
\begin{equation}\label{eq:new-rule}
\begin{aligned}
&m(v)>B_A(v;m)\quad\Longrightarrow\quad m(w)\ge B_A(w;m)+g(A,k),\\
&g(A,k)=\begin{cases}2k+1,&A\text{ self-reversed},\\ k+1,&\text{otherwise}.
\end{cases}
\end{aligned}
\end{equation}
Here self-reversed means that for some integer $c$,
$\cm(A(j))=A(c-1-j)$ for every $j$; equivalently its reverse complement
is a translate of itself. Take the maximum of all bounds for both words
of an orbit and assign it to both. At a fixed word round upward to even.
Every bound refers only to longer words, so this defines an explicit
recursion. Enumerate every $v\in V$, $1\le p<|v|$, and $k$ with
$kp<|v|$, retaining a rule precisely when repetition of $v[0:p)$ agrees
with $v$. This includes all nonleast periods and all phases represented
by such words. No fixed-word-only $2k+2$ variant is used.

\begin{lemma}[Tagged-text domination]\label{lem:domination}
For every feasible $T$ and every nonempty $s\in V$,
\[
m(s)\le N(s):=N_T(s)+N_{\rho(T)}(s).
\]
Thus $W:=\sum_{a\in\Sigma}m(a)\le2\OPT(\mathcal S)$, and $m(s)\le2L$.
\end{lemma}
\begin{proof}
Treat $T$ and $\rho(T)$ as distinct \emph{tagged components}, even if
their word labels coincide. Then $N$ is symmetric and even at every
fixed word. It satisfies the required and extension bounds. Descending
induction for the ordinary $k+1$ rule is the proof of
\src{\texttt{lem:counts}} on this disjoint collection: an unblocked
$v$ gives $k+1$ period-spaced unblocked prefixes, or a positive
longer-bracket deficit contributes that many placements. The blocking
functional is linear over the components.

We prove the additional $2k+1$ bound. Translate so that $v=A[0:h)$,
$|w|=h-kp$. In a self-reversed text, $\rho(v)$ occupies $A[c-h:c)$.
The possible starts of a copy of $w$ inside this interval form an
integer interval of $kp+1$ integers. At least $k$ of them have residue
$0$ modulo $p$, and hence spell $w$ in the required $A$ alignment.
The original $v$ contains $k+1$ such copies, at $0,p,\ldots,kp$.

If $N(v)>B_A(v;N)$, choose an unblocked $v$ in one tagged component.
All $k+1$ internal copies are unblocked: a mismatch bracket of an
internal copy must enclose the entire matching $v$ interval. Reflect
this occurrence into the other component. Each of the indicated $k$
aligned copies in $\rho(v)$ is likewise unblocked. A bracket blocking
one would enclose $\rho(v)$; its reflection is in $V$ and would block
the original $v$. End mismatches remain mismatches because $\cm$ is a
bijection. Tagging makes the two sets of occurrences disjoint. Therefore
\[
N(w)-B_A(w;N)\ge(k+1)+k=2k+1.
\]

Otherwise $N(v)=B_A(v;N)\ge m(v)>B_A(v;m)$. Write
$\Delta(r)=N(r)-m(r)$ for longer words. Induction makes these
nonnegative. Some bracket word $r$ has $b_{A,v}(r)>0$ and
$\Delta(r)>0$. Its chosen placement contains $k+1$ counted placements
of $w$. The reflected bracket in $\rho(r)$ contains at least $k$
of the aligned shifted-phase placements just described; both endpoints
remain mismatches, so they are counted by the same functional for $w$.
If $r\ne\rho(r)$, symmetry gives
$\Delta(\rho(r))=\Delta(r)\ge1$, and these two bracket terms contribute
at least $2k+1$ to $B_A(w;N)-B_A(w;m)$. If $r=\rho(r)$, the positive
integer $\Delta(r)$ is even, hence at least $2$; the original $k+1$
placements alone give $2(k+1)\ge2k+1$. There is no double counting of
a fixed bracket. All other terms are nonnegative. This proves the new
bound in both cases.

Orbit maxima stay below the symmetric $N$, and even rounding stays below
its even values at fixed words. Induction follows. Summing letter
counts gives $W\le2|T|$. Concatenating one surviving representative per
orbit supplies a feasible word of length $L$, proving $m(s)\le2L$.
\end{proof}

\subsection{The balanced graph and fixed vertices}\label{sec:base}
For each nonempty $s$, the hierarchical graph has an up edge
$\pref(s)\to s$ and a down edge $s\to\suf(s)$; prefix and suffix
here delete the last and first letter respectively. Up cost is one,
down cost zero. Give these edges multiplicities
\begin{equation}\label{eq:multiplicities}
u(s)=m(s)-\sum_c m(cs),\qquad d(s)=m(s)-\sum_c m(sc).
\end{equation}
By the telescoping proof of \src{\texttt{lem:base}}, this graph $G$ is
balanced, has $W$ up and $W$ down copies, and uses every required vertex.
Indeed both degrees at nonempty $s$ equal
$m(s)-\sum_{c,d}m(csd)$, and the total multiplicities telescope to $W$.

Symmetry gives $u(s)=d(\rho(s))$. Pair copies to define an
edge-reversing involution $\tau$, mapping the up edge at $s$ to the
down edge at $\rho(s)$. There are no fixed edges. At a nonempty fixed
$s$, the displayed degree is even: nonfixed extension words pair and
fixed extension words have even counts. At the empty vertex,
\begin{equation}\label{eq:empty-degree}
\outdeg_G(\eps)=\sum_a u(a)
=\sum_{|s|=1}m(s)-\sum_{|s|=2}m(s)
\end{equation}
is also even. It is not the total up cost $W$. Adding a closed walk
$C$ together with its reversed image $\tau C$ preserves balance,
symmetry, and even outdegree at fixed vertices, since the increment
there is $\outdeg_C+\indeg_C=2\outdeg_C$.
